\documentclass[11pt]{article}
\usepackage[a4paper,margin=1in]{geometry}
\usepackage{amsmath,amssymb,mathtools}
\usepackage{enumitem}
\usepackage{unicode-math}
\setlist{leftmargin=*}
\allowdisplaybreaks
\emergencystretch=3em
\title{Sharp KKT Bound in the Small-Parameter Square}
\author{}
\date{}
\begin{document}
\maketitle
I pushed the Sonin/Bessel route one step further. The useful output is a \textbf{closed proof of the sharp KKT refined bound in the small-parameter square}
\begin{equation*}
0\le \alpha,\beta\le \frac12,\qquad n\ge1.
\end{equation*}
This is not the full conjecture, but it is a genuine analytic subcase using the direct \(g\)-Sonin functional plus endpoint Bessel comparison. It avoids integer-parameter representation theory. The target object and refined bound are the ones summarized in your KKT review.

\section*{Theorem proved here}

Let
\begin{equation*}
g(x)=g_n^{(\alpha,\beta)}(x)
=
C_{n,\alpha,\beta}
\left(\frac{1-x}{2}\right)^{\alpha/2}
\left(\frac{1+x}{2}\right)^{\beta/2}
P_n^{(\alpha,\beta)}(x),
\end{equation*}
where
\begin{equation*}
C_{n,\alpha,\beta}^2
=
\frac{\Gamma(n+1)\Gamma(n+\alpha+\beta+1)}
{\Gamma(n+\alpha+1)\Gamma(n+\beta+1)}.
\end{equation*}
Then for
\begin{equation*}
n\ge1,\qquad 0\le\alpha,\beta\le\frac12,
\end{equation*}
one has
\begin{equation*}
\boxed{
|g_n^{(\alpha,\beta)}(x)|
\le
\left(
\frac{(n+1)(n+\alpha+\beta+1)}
{(n+\alpha+1)(n+\beta+1)}
\right)^{1/4}
}
\qquad (-1\le x\le1).
\end{equation*}
This is the refined KKT constant, not merely the weaker \(|g|\le1\).

\section*{1. Direct Sturm--Liouville equation for \(g\)}

Set
\begin{equation*}
s=\alpha+\beta,\qquad d=\beta-\alpha,
\end{equation*}
\begin{equation*}
A(x)=1-x^2,
\end{equation*}
\begin{equation*}
N=n+\frac{\alpha+\beta+1}{2},
\end{equation*}
and
\begin{equation*}
U(x)=d-sx=\beta-\alpha-(\alpha+\beta)x.
\end{equation*}

Then \(g\) satisfies
\begin{equation}
\boxed{
(A g')'+G(x)g=0
}
\tag{1}
\end{equation}
where
\begin{equation}
\boxed{
G(x)
=
N^2-\frac14
-\frac{\alpha^2}{2(1-x)}
-\frac{\beta^2}{2(1+x)}
}
\tag{2}
\end{equation}

Equivalently,
\begin{equation*}
G(x)=\kappa-\frac{U(x)^2}{4A(x)},
\end{equation*}
with
\begin{equation*}
\kappa=n(n+\alpha+\beta+1)+\frac{\alpha+\beta}{2}.
\end{equation*}

Define
\begin{equation}
\boxed{
F(x)=A(x)G(x)
=
\kappa(1-x^2)-\frac14\bigl(d-sx\bigr)^2
}
\tag{3}
\end{equation}

The direct turning points are the roots of \(F\):
\begin{equation}
\boxed{
x_\pm
=
\frac{ds\pm4\sqrt{\kappa(\kappa+\alpha\beta)}}{s^2+4\kappa}
}
\tag{4}
\end{equation}

For \(\alpha,\beta>0\),
\begin{equation*}
F(-1)=-\beta^2,\qquad F(1)=-\alpha^2,
\end{equation*}
so
\begin{equation*}
F(x)>0
\quad\Longleftrightarrow\quad
x\in(x_-,x_+).
\end{equation*}

The vertex is
\begin{equation}
\boxed{
x_0=\frac{ds}{s^2+4\kappa}
}
\tag{5}
\end{equation}

\section*{2. Correct Sonin functional for \(g\)}

On \(F>0\), define
\begin{equation}
\boxed{
S_g(x)
=
g(x)^2+\frac{(1-x^2)^2g'(x)^2}{F(x)}
}
\tag{6}
\end{equation}

A direct differentiation using (1) gives
\begin{equation}
\boxed{
S_g'(x)
=
-\frac{(1-x^2)^2F'(x)}{F(x)^2}\,g'(x)^2
}
\tag{7}
\end{equation}

Since \(F\) is a concave quadratic,
\begin{equation*}
F'(x)>0\quad\text{on }(x_-,x_0),
\end{equation*}
and
\begin{equation*}
F'(x)<0\quad\text{on }(x_0,x_+).
\end{equation*}

Therefore
\begin{equation}
\boxed{
S_g \text{ decreases on }(x_-,x_0),
\qquad
S_g \text{ increases on }(x_0,x_+).
}
\tag{8}
\end{equation}

At a critical point \(g'(x_*)=0\),
\begin{equation*}
S_g(x_*)=g(x_*)^2.
\end{equation*}

Hence the largest local extrema on the right side occur in the \textbf{first endpoint lobe near \(x=1\)}, and the largest local extrema on the left side occur in the \textbf{first endpoint lobe near \(x=-1\)}.

So the problem reduces to bounding the two endpoint lobes.

\section*{3. Right-endpoint Bessel comparison}

Put
\begin{equation*}
x=\cos\theta,\qquad 0<\theta<\pi,
\end{equation*}
and define
\begin{equation*}
v(\theta)=\sqrt{\sin\theta}\,g(\cos\theta).
\end{equation*}

Then
\begin{equation}
\boxed{
v''(\theta)+Q(\theta)v(\theta)=0
}
\tag{9}
\end{equation}
with
\begin{equation}
\boxed{
Q(\theta)
=
N^2
+
\frac{1/4-\alpha^2}{4\sin^2(\theta/2)}
+
\frac{1/4-\beta^2}{4\cos^2(\theta/2)}
}
\tag{10}
\end{equation}

Assume now
\begin{equation*}
0\le \alpha,\beta\le\frac12.
\end{equation*}

Then
\begin{equation*}
1/4-\alpha^2\ge0,\qquad 1/4-\beta^2\ge0.
\end{equation*}

Also,
\begin{equation*}
4\sin^2(\theta/2)\le \theta^2.
\end{equation*}

Therefore
\begin{equation}
Q(\theta)
\ge
N^2+\frac{1/4-\alpha^2}{\theta^2}
=:q_\alpha(\theta).
\tag{11}
\end{equation}

The comparison equation
\begin{equation*}
z''+q_\alpha(\theta)z=0
\end{equation*}
has the regular Bessel solution
\begin{equation*}
z(\theta)=K_\alpha\sqrt{\theta}\,J_\alpha(N\theta).
\end{equation*}

The constant matching the leading behavior of \(v\) at \(\theta=0\) is
\begin{equation}
\boxed{
K_\alpha
=
C_{n,\alpha,\beta}
\frac{\Gamma(n+\alpha+1)}
{\Gamma(n+1)N^\alpha}
}
\tag{12}
\end{equation}

Indeed, as \(\theta\downarrow0\),
\begin{equation*}
v(\theta)
\sim
C_{n,\alpha,\beta}
2^{-\alpha}
\frac{\Gamma(n+\alpha+1)}
{\Gamma(\alpha+1)\Gamma(n+1)}
\theta^{\alpha+1/2},
\end{equation*}
while
\begin{equation*}
K_\alpha\sqrt{\theta}J_\alpha(N\theta)
\sim
K_\alpha
\frac{N^\alpha}{2^\alpha\Gamma(\alpha+1)}
\theta^{\alpha+1/2}.
\end{equation*}

By the Sturm comparison/Wronskian argument, while both \(v\) and \(z\) are positive,
\begin{equation*}
\frac{v}{z}
\end{equation*}
is decreasing from \(1\). Therefore, throughout the first right endpoint lobe,
\begin{equation}
\boxed{
0\le g(\cos\theta)
\le
K_\alpha
\left(\frac{\theta}{\sin\theta}\right)^{1/2}
J_\alpha(N\theta)
}
\tag{13}
\end{equation}

So the right endpoint lobe is controlled by a scalar Bessel envelope.

\section*{4. Bounding the Bessel envelope}

Let \(0\le a\le1/2\). For \(0\le t\le j_{a,1}\), the first positive zero of \(J_a\), the normalized Bessel function satisfies
\begin{equation}
\boxed{
J_a(t)
\le
\frac{(t/2)^a}{\Gamma(a+1)}
\exp\!\left(-\frac{t^2}{4(a+1)}\right).
}
\tag{14}
\end{equation}

This follows from the Hadamard product
\begin{equation*}
2^a\Gamma(a+1)t^{-a}J_a(t)
=
\prod_{k=1}^{\infty}
\left(1-\frac{t^2}{j_{a,k}^2}\right),
\end{equation*}
the Rayleigh sum
\begin{equation*}
\sum_{k=1}^{\infty}\frac1{j_{a,k}^2}
=
\frac1{4(a+1)},
\end{equation*}
and \(1-y\le e^{-y}\).

Since \(0\le a\le1/2\), one has
\begin{equation*}
j_{a,1}\le j_{1/2,1}=\pi.
\end{equation*}
Also \(n\ge1\), so
\begin{equation*}
N=n+\frac{\alpha+\beta+1}{2}\ge\frac32.
\end{equation*}
Thus in the first endpoint lobe,
\begin{equation*}
0\le \theta\le \frac{j_{\alpha,1}}{N}\le\frac{2\pi}{3}.
\end{equation*}

On \(0\le\theta\le2\pi/3\), the elementary estimate
\begin{equation}
\boxed{
\left(\frac{\theta}{\sin\theta}\right)^{1/2}
\le
e^{\theta^2/8}
}
\tag{15}
\end{equation}
holds.

Using (14) and (15) in (13), with \(a=\alpha\), gives
\begin{equation}
g(\cos\theta)
\le
D_{n,\alpha,\beta}\,
\theta^\alpha
\exp\!\left(
-\left[
\frac{N^2}{4(\alpha+1)}-\frac18
\right]\theta^2
\right),
\tag{16}
\end{equation}
where
\begin{equation*}
D_{n,\alpha,\beta}
=
C_{n,\alpha,\beta}
\frac{\Gamma(n+\alpha+1)}
{\Gamma(n+1)2^\alpha\Gamma(\alpha+1)}.
\end{equation*}

Let
\begin{equation*}
c=
\frac{N^2}{4(\alpha+1)}-\frac18
=
\frac{2N^2-\alpha-1}{8(\alpha+1)}.
\end{equation*}
Then \(c>0\), and
\begin{equation*}
\max_{\theta\ge0}\theta^\alpha e^{-c\theta^2}
=
\left(\frac{\alpha}{2ec}\right)^{\alpha/2},
\end{equation*}
with the value interpreted as \(1\) when \(\alpha=0\).

Therefore
\begin{equation}
g(\cos\theta)
\le
D_{n,\alpha,\beta}
\left(
\frac{4\alpha(\alpha+1)}
{e(2N^2-\alpha-1)}
\right)^{\alpha/2}.
\tag{17}
\end{equation}

\section*{5. Gamma-ratio estimate gives the exact KKT constant}

We now show that the right-hand side of (17) is bounded by the KKT constant.

First,
\begin{equation*}
D_{n,\alpha,\beta}
=
\frac{1}{2^\alpha\Gamma(\alpha+1)}
\left[
\frac{\Gamma(n+\alpha+\beta+1)\Gamma(n+\alpha+1)}
{\Gamma(n+\beta+1)\Gamma(n+1)}
\right]^{1/2}.
\end{equation*}

By Wendel's gamma-ratio inequality, for \(0\le\alpha\le1\),
\begin{equation*}
\frac{\Gamma(n+\alpha+1)}{\Gamma(n+1)}
\le
(n+1)^\alpha,
\end{equation*}
and
\begin{equation*}
\frac{\Gamma(n+\alpha+\beta+1)}{\Gamma(n+\beta+1)}
\le
(n+\beta+1)^\alpha.
\end{equation*}

Hence
\begin{equation*}
D_{n,\alpha,\beta}
\le
\frac{\bigl((n+1)(n+\beta+1)\bigr)^{\alpha/2}}
{2^\alpha\Gamma(\alpha+1)}.
\end{equation*}

Thus the right endpoint envelope is bounded by
\begin{equation}
B_R
=
\frac{1}{\Gamma(\alpha+1)}
\left[
\frac{\alpha(\alpha+1)(n+1)(n+\beta+1)}
{e(2N^2-\alpha-1)}
\right]^{\alpha/2}.
\tag{18}
\end{equation}

Now use the elementary inequality
\begin{equation}
\boxed{
\frac{\alpha(\alpha+1)(n+1)(n+\beta+1)}
{e(2N^2-\alpha-1)}
\le
\frac{\alpha}{2}
}
\tag{19}
\end{equation}
for
\begin{equation*}
n\ge1,\qquad 0\le\alpha,\beta\le\frac12.
\end{equation*}

One way to verify (19) is to use \(e>8/3\) and check that
\begin{equation*}
2N^2-\alpha-1
\ge
\frac34(\alpha+1)(n+1)(n+\beta+1).
\end{equation*}
Writing \(m=n+1\ge2\), the difference equals
\begin{align*}
&2N^2-\alpha-1
-\frac34(\alpha+1)m(m+\beta) \\
&\quad =
\left(\frac54-\frac{3\alpha}{4}\right)(m-2)^2
+
\left(3-\alpha+\frac{5\beta}{4}-\frac{3\alpha\beta}{4}\right)(m-2) \\
&\qquad
+
\frac{(1-\alpha)^2}{2}
+
\frac{\beta^2}{2}
+
\frac{(3-\alpha)\beta}{2}.
\end{align*}
Every term is nonnegative for \(m\ge2\) and \(0\le\alpha,\beta\le1/2\).

Therefore
\begin{equation}
B_R
\le
\frac{1}{\Gamma(\alpha+1)}
\left(\frac{\alpha}{2}\right)^{\alpha/2}.
\tag{20}
\end{equation}

Using convexity of \(\log\Gamma\),
\begin{equation*}
\log\Gamma(1+\alpha)\ge -\gamma\alpha,
\end{equation*}
where \(\gamma\) is Euler's constant. Hence
\begin{equation*}
\log B_R
\le
\frac{\alpha}{2}\log\frac{\alpha}{2}
+\gamma\alpha.
\end{equation*}
Since \(\alpha\le1/2\),
\begin{equation*}
\frac12\log\frac{\alpha}{2}
\le
-\log2.
\end{equation*}
Thus
\begin{equation}
\log B_R
\le
-(\log2-\gamma)\alpha.
\tag{21}
\end{equation}

Now define the KKT refined constant
\begin{equation*}
T_{n,\alpha,\beta}
=
\left(
\frac{(n+1)(n+\alpha+\beta+1)}
{(n+\alpha+1)(n+\beta+1)}
\right)^{1/4}.
\end{equation*}

Observe that
\begin{equation*}
T_{n,\alpha,\beta}^4
=
1-
\frac{\alpha\beta}{(n+\alpha+1)(n+\beta+1)}.
\end{equation*}
Let
\begin{equation*}
z=
\frac{\alpha\beta}{(n+\alpha+1)(n+\beta+1)}.
\end{equation*}
Since \(n\ge1\), \(0\le\beta\le1/2\), and \(0\le\alpha\le1/2\),
\begin{equation*}
0\le z\le \frac{\alpha}{8}\le\frac1{16}.
\end{equation*}
Therefore
\begin{equation}
\log T_{n,\alpha,\beta}
=
\frac14\log(1-z)
\ge
-\frac{z}{4(1-z)}
\ge
-\frac{\alpha}{30}.
\tag{22}
\end{equation}

Finally,
\begin{equation*}
\log2-\gamma>\frac1{30}.
\end{equation*}
Combining (21) and (22),
\begin{equation*}
\log B_R
\le
-(\log2-\gamma)\alpha
\le
-\frac{\alpha}{30}
\le
\log T_{n,\alpha,\beta}.
\end{equation*}
Hence
\begin{equation}
\boxed{
B_R\le T_{n,\alpha,\beta}.
}
\tag{23}
\end{equation}

So the first right endpoint lobe satisfies the sharp KKT refined bound.

\section*{6. Left endpoint}

The left endpoint is identical after swapping
\begin{equation*}
\alpha\leftrightarrow\beta,
\qquad
x\leftrightarrow -x.
\end{equation*}

Thus the first left endpoint lobe is also bounded by
\begin{equation*}
T_{n,\alpha,\beta}.
\end{equation*}

By the direct \(g\)-Sonin localization, every interior local extremum is bounded by one of these two endpoint-lobe bounds. Endpoint values are harmless:

\begin{itemize}
\item if \(\alpha,\beta>0\), then \(g(\pm1)=0\);
\item if \(\alpha=0\), then \(g(1)=1\) and the refined constant is \(1\);
\item if \(\beta=0\), then \(g(-1)=(-1)^n\) and the refined constant is \(1\).
\end{itemize}

Therefore
\begin{equation*}
\boxed{
|g_n^{(\alpha,\beta)}(x)|
\le
T_{n,\alpha,\beta}
}
\end{equation*}
for all \(x\in[-1,1]\), \(n\ge1\), and
\begin{equation*}
0\le\alpha,\beta\le\frac12.
\end{equation*}

This proves the refined KKT inequality in the full small-parameter square.

\section*{What this adds to the program}

The A2 attempts had the correct angular Liouville equation but used the wrong function globally. The corrected strategy is:

\begin{enumerate}
\item Use the \textbf{direct \(g\)-equation} and the \(F\)-Sonin functional to localize the maximum to endpoint lobes.
\item Use the \textbf{Liouville-transformed equation} only inside those endpoint lobes.
\item Compare with a regular Bessel model.
\item Use gamma-ratio estimates to recover the exact KKT constant.
\end{enumerate}

This now gives two complementary proven regimes:
\begin{equation*}
\boxed{
0\le\alpha,\beta\le\frac12
}
\end{equation*}
by the endpoint Bessel comparison above, and
\begin{equation*}
\boxed{
\alpha,\beta\ge2n
}
\end{equation*}
by the earlier weighted-energy argument.

The remaining hard region is therefore the middle/asymmetric zone:
\begin{equation*}
\max(\alpha,\beta)>\frac12
\quad\text{and}\quad
\min(\alpha,\beta)<2n,
\end{equation*}
especially when one parameter is small and the other is large.

The next technical target should be to replace the simple comparison
\begin{equation*}
Q(\theta)\ge N^2+\frac{1/4-\alpha^2}{\theta^2}
\end{equation*}
with a \textbf{shifted Bessel comparison} that tolerates the negative finite endpoint contribution
\begin{equation*}
\frac{1/4-\beta^2}{4\cos^2(\theta/2)}
\end{equation*}
when \(\beta>1/2\). That is the first obstruction outside the square \(0\le\alpha,\beta\le1/2\).
\end{document}
